\documentclass[12pt]{article}

\usepackage{sbc-template}
\usepackage{graphicx,url}
\usepackage[utf8]{inputenc}
\usepackage[brazil]{babel}
\usepackage{amsmath,amsfonts,amssymb}
\usepackage{booktabs}
\usepackage{subcaption}
\usepackage{hyperref}
\usepackage{cite}
\usepackage{multirow}
\usepackage{array}

\hypersetup{
    colorlinks=true,
    linkcolor=blue,
    citecolor=blue,
    urlcolor=blue
}

\title{Reprodução Experimental e Avaliação de Mapas de Ativação Desacoplados para Detecção de Massas em Mamografia via Aprendizado Fracamente Supervisionado}

\author{Grupo 4\inst{1}}

\address{Campus Senador Helvídio Nunes de Barros (CSHNB)\\
  Universidade Federal do Piauí (UFPI) -- Picos, PI -- Brasil\\
  Disciplina: Tópicos Especiais em Visão Computacional\\
  Docente: Prof. Dr. José Denes Lima Araújo
  \email{\{grupo4.visao\}@ufpi.edu.br}
}

\begin{document}

\maketitle

\begin{abstract}
A detecção precoce do câncer de mama por Mamografia Digital Direta (FFDM) é primordial para a redução da mortalidade feminina. Contudo, modelos de detecção de objetos totalmente supervisionados esbarram no custo proibitivo da anotação de contornos por radiologistas. O Aprendizado Fracamente Supervisionado para Detecção de Objetos (WSOD) surge como solução elegante, treinando redes unicamente com o rótulo binário da imagem. Este trabalho apresenta a reprodução científica e a análise crítica experimental do artigo de Sampaio e Cordeiro (UFRPE, SBC 2022), o qual estende o modelo GMIC (\textit{Globally-Aware Multiple Instance Classifier}) e introduz a hipótese do desacoplamento de Mapas de Ativação de Classe (CAM) entre as fases de treinamento e teste. Implementou-se uma arquitetura modular em PyTorch integrando cinco métodos de saliência convolucional (CAM, Grad-CAM, Grad-CAM++, XGrad-CAM e Layer-CAM) avaliados no benchmark internacional VinDr-Mammo. Para contornar restrições de memória de vídeo em imagens de alta resolução ($2944 \times 1920$), empregou-se Precisão Mista Automática (AMP FP16) e Acumulação de Gradientes. Os resultados confirmam empiricamente a superioridade do modelo campeão desacoplado GMIC(XGradCAM) com teste em Grad-CAM++, alcançando Sensibilidade de 0,70 com uma redução de 43,2\% nos falsos positivos (0,88 FPPI contra 1,55 do baseline). Adicionalmente, apresenta-se uma demonstração analítica corrigindo uma inconsistência técnica na Tabela 1 do artigo original relacionada à taxa de falsos negativos.
\end{abstract}

\begin{abstracten}
Early detection of breast cancer through Full-Field Digital Mammography (FFDM) is essential to reduce female cancer mortality. However, fully supervised object detection frameworks face severe clinical annotation bottlenecks. Weakly Supervised Object Detection (WSOD) overcomes this limitation by training deep neural models using exclusively image-level classification labels. This paper presents an experimental reproduction and critical analysis of the study by Sampaio and Cordeiro (UFRPE, SBC 2022), which investigates the GMIC (Globally-Aware Multiple Instance Classifier) framework and introduces the hypothesis of decoupling Class Activation Mapping (CAM) methods between training and inference phases. We developed a modular PyTorch architecture evaluating five convolutional saliency techniques (CAM, Grad-CAM, Grad-CAM++, XGrad-CAM, and Layer-CAM) on the VinDr-Mammo benchmark. To mitigate GPU memory constraints when processing ultra-high-resolution mammograms ($2944 \times 1920$), Automatic Mixed Precision (AMP FP16) and Gradient Accumulation were implemented. The experimental findings confirm the superiority of the decoupled champion model GMIC(XGradCAM) evaluated with Grad-CAM++, achieving a True Positive Rate (TPR) of 0.70 with a 43.2\% reduction in false alarms (0.88 FPPI versus 1.55 from the baseline). Furthermore, an analytical proof is presented correcting a technical errata in Table 1 of the original publication concerning the reported false negative rate.
\end{abstracten}

\noindent\textbf{Palavras-chave:} Aprendizado Fracamente Supervisionado, Detecção de Massas Mamárias, GMIC, Mapas de Ativação de Classe, Reprodutibilidade Científica.

% =========================================================================
% SEÇÃO 1: INTRODUÇÃO
% =========================================================================
\section{Introdução}

O câncer de mama permanece como a neoplasia maligna de maior incidência e a principal causa de mortalidade oncológica entre as mulheres globalmente, respondendo por aproximadamente 15,5\% do total de óbitos femininos decorrentes de neoplasias malignas \cite{who2023breast}. Estudos epidemiológicos consolidados demonstram que a sobrevida em cinco anos atinge patamares superiores a 90\% quando o tumor é diagnosticado em fases iniciais (lesões inferiores a 10 mm), ao passo que a disseminação metastática tardia reduz essa perspectiva para índices inferiores a 30\%.

A Mamografia Digital Direta (\textit{Full-Field Digital Mammography} - FFDM) consolidou-se internacionalmente como a modalidade padrão-ouro nos programas de rastreamento populacional (\textit{screening}). Todavia, a interpretação radiológica de mamogramas digitais é notória por sua complexidade: o fenômeno de mascaramento gerado pela sobreposição de tecido fibroglandular denso, associado à sutileza morfológica de lesões iniciais com contornos espiculados ou microlobulados, eleva substancialmente a fadiga cognitiva dos médicos especialistas. Estima-se que até 30\% das lesões malignas passem despercebidas durante sessões rotineiras de triagem visual \cite{ribli2018detecting}. Nesse cenário, sistemas de Diagnóstico Auxiliado por Computador (\textit{Computer-Aided Diagnosis} - CAD) baseados em Visão Computacional atuam como uma segunda opinião automatizada, delimitando regiões suspeitas e auxiliando na redução de omissões diagnósticas.

Não obstante o avanço das Redes Neurais Convolucionais (CNNs), a implementação tradicional de modelos de detecção de objetos (\textit{Fully Supervised Object Detection}) enfrenta o denominado ``gargalo da anotação médica'' \cite{shen2021interpretable}. A delimitação manual de caixas delimitadoras (\textit{bounding boxes}) ou máscaras de segmentação pixel a pixel em imagens médicas de altíssima resolução ($2944 \times 1920$ pixels) requer milhares de horas de radiologistas seniores certificados, gerando custos financeiros proibitivos e expressiva variabilidade interobservador. Como consequência, menos de 1\% das bases de dados radiológicas arquivadas em hospitais dispõe de coordenadas geométricas de anotação, embora existam milhões de exames laudados no nível do exame.

Para superar essa barreira, o paradigma de Aprendizado Fracamente Supervisionado para Detecção de Objetos (\textit{Weakly Supervised Object Detection} - WSOD) propõe treinar modelos profundos utilizando estritamente o rótulo binário global da imagem ($y=0$: exame normal; $y=1$: presença de massa tumoral), sem fornecer nenhuma informação espacial de localização durante o treinamento. A partir de mecanismos de explicabilidade e saliência convolucional, o modelo é incumbido de inferir a localização espacial exata da lesão na fase de teste.

O presente trabalho dedica-se à reprodução científica rigorosa, auditoria técnica e ampliação metodológica do artigo intitulado \textit{``A Study on Class Activation Map Methods to Detect Masses in Mammography Images using Weakly Supervised Learning''}, de autoria de Vicente Sampaio e Filipe R. Cordeiro (UFRPE), publicado nos anais da Sociedade Brasileira de Computação (SBC) \cite{sampaio2022study}. O estudo em questão investiga a arquitetura GMIC (\textit{Globally-Aware Multiple Instance Classifier}), concebida por Shen et al. \cite{shen2021interpretable}, avaliando o comportamento de cinco técnicas modernas de mapas de ativação de classe (CAM) no benchmark internacional VinDr-Mammo \cite{vindrmammo2022}.

A hipótese científica fulcral investigada consiste no \textbf{desacoplamento de CAMs}: demonstrar que o método de ativação convolucional ideal para selecionar recortes de alta resolução durante o treinamento difere funcionalmente daquele ideal para delimitar a caixa delimitadora final na inferência. Este artigo reconstrói a base de código integralmente em PyTorch, introduz adaptações de engenharia para viabilizar o processamento de imagens de ultra-resolução em GPUs convencionais via Precisão Mista Automática (AMP) e Acumulação de Gradientes, valida a matriz mestra de resultados e formaliza uma errata matemática relativa à métrica de falsos negativos da publicação original.

% =========================================================================
% SEÇÃO 2: FUNDAMENTAÇÃO TEÓRICA
% =========================================================================
\section{Fundamentação Teórica e Trabalhos Relacionados}

\subsection{Formulação Matemática do WSOD e Multiple Instance Learning}

No regime WSOD para diagnóstico médico, o conjunto de dados de treino é expresso por $\mathcal{D}_{train} = \{(x_n, y_n)\}_{n=1}^N$, onde $x_n \in \mathbb{R}^{H \times W \times 1}$ representa a imagem mamográfica em escala de cinza e $y_n \in \{0, 1\}$ é o rótulo supervisionado global. Durante o treino, o vetor de coordenadas de qualquer lesão $\mathcal{B}^* = \{[x_{min}, y_{min}, x_{max}, y_{max}]\}$ permanece oculto à função de perda.

O problema é comumente formulado sob o arcabouço de Aprendizado por Múltiplas Instâncias (\textit{Multiple Instance Learning} - MIL) \cite{ilse2018attention}. Cada mamograma $x_n$ é modelado como uma ``sacola'' (\textit{bag}) composta por múltiplas instâncias espaciais locais $\{p_k\}_{k=1}^K$. A hipótese canônica de MIL estabelece que a sacola é rotulada como positiva ($y=1$) se, e somente se, pelo menos uma de suas instâncias contiver tecido tumoral maligno:
\begin{equation}
y = \max_{k \in \{1, \dots, K\}} y^{(k)}
\end{equation}

\subsection{A Arquitetura GMIC}

O modelo GMIC (\textit{Globally-Aware Multiple Instance Classifier}) foi concebido por Shen et al. \cite{shen2021interpretable} com o propósito de harmonizar duas exigências conflitantes no diagnóstico mamográfico: capturar o contexto global da glândula mamária sem descartar os detalhes ultra-resolvidos de lesões milimétricas, mantendo o consumo de memória de vídeo computacionalmente viável. A arquitetura organiza-se em três módulos articulados:

\begin{enumerate}
    \item \textbf{Módulo Global}: Composto por uma ResNet-22 personalizada para imagens de alta resolução, recebe o exame completo $x \in \mathbb{R}^{H \times W}$ e produz um mapa de características convolucionais $S_g \in \mathbb{R}^{C_g \times H_g \times W_g}$ com $C_g = 512$. Após uma camada de \textit{Global Average Pooling} (GAP), gera a representação contextual $h_{global} \in \mathbb{R}^{512}$ e a predição global $\hat{y}_{global} = \sigma(W_g h_{global} + b_g)$.
    
    \item \textbf{Extrator Guloso de Patches}: A partir do mapa de características globais, projeta-se o Mapa de Ativação de Classe $A^c \in \mathbb{R}^{H_g \times W_g}$ correspondente à classe massa neoplásica. Um algoritmo guloso (\textit{greedy search}) com supressão de vizinhança identifica as $K=6$ coordenadas de pico térmico $(c_y, c_x)$ e recorta $K$ patches locais de dimensão $256 \times 256$ pixels diretamente da imagem original na resolução nativa.
    
    \item \textbf{Módulo Local e Fusão}: Uma ResNet-18 independente processa individualmente cada um dos $K$ patches recortados, extraindo representações convolucionais $h_p \in \mathbb{R}^{512}$ que capturam margens espiculadas e microcalcificações pleomórficas. A agregação MIL é realizada via operador Max-Pooling sobre as instâncias:
    \begin{equation}
    h_{local} = \max_{k \in \{1, \dots, K\}} h_{p,k}
    \end{equation}
    O Módulo de Fusão concatena os vetores $[h_{global}; h_{local}] \in \mathbb{R}^{1024}$ e os submete a camadas densas seguidas de ativação Sigmoid para produzir a predição consolidada $\hat{y}_{fusao}$.
\end{enumerate}

\subsection{Função de Perda Multi-Ramo e Regularização de Esparsidade L1}

Para compelir todos os módulos a convergirem de maneira sinérgica, o treinamento otimiza conjuntamente as predições dos três ramos através da Entropia Cruzada Binária ($\text{BCE}$). Para prevenir que o mapa de ativação $A^c$ se expanda indiscriminadamente por todo o parênquima sadio, adiciona-se uma penalização de esparsidade baseada na norma L1 ponderada pelo hiperparâmetro $\beta = 3,26$:
\begin{equation}
\mathcal{L}(y, \hat{y}) = \text{BCE}(y, \hat{y}_{local}) + \text{BCE}(y, \hat{y}_{global}) + \text{BCE}(y, \hat{y}_{fusao}) + \beta \mathcal{L}_{reg}(A^c)
\end{equation}
onde:
\begin{equation}
\mathcal{L}_{reg}(A^c) = \frac{1}{H_g W_g} \sum_{i=1}^{H_g} \sum_{j=1}^{W_g} |A_{i,j}^c|
\end{equation}
Caso $\beta = 0$, a rede global minimizaria a perda de classificação ativando áreas dispersas de tecido normal, corrompendo a pureza dos patches locais e inviabilizando a convergência da ResNet-18.

\subsection{Taxonomia dos Cinco Métodos de Mapas de Ativação (CAM)}

A contribuição fulcral de Sampaio e Cordeiro \cite{sampaio2022study} reside no estudo comparativo de cinco técnicas de interpretabilidade visual baseadas em ativação e gradiente:

\begin{enumerate}
    \item \textbf{CAM Clássico \cite{zhou2016learning}}: Requer obrigatoriamente que a camada linear final seja conectada diretamente ao GAP dos mapas convolucionais:
    \begin{equation}
    M_{\text{CAM}}^c(x, y) = \text{ReLU}\left( \sum_k w_k^c A^k(x, y) \right)
    \end{equation}
    Gera ativações excessivamente difusas e volumosas em imagens mamográficas.

    \item \textbf{Grad-CAM \cite{selvaraju2017grad}}: Generaliza o CAM para qualquer topologia convolucional através da média espacial dos gradientes de primeira ordem:
    \begin{equation}
    \alpha_k^c = \frac{1}{Z} \sum_{i} \sum_{j} \frac{\partial Y^c}{\partial A_{i,j}^k}, \quad M_{\text{Grad-CAM}}^c = \text{ReLU}\left( \sum_k \alpha_k^c A^k \right)
    \end{equation}
    Tende a focar exclusivamente no núcleo central de maior contraste, subestimando margens espiculadas.

    \item \textbf{Grad-CAM++ \cite{chattopadhay2018grad}}: Utiliza derivadas parciais analíticas de segunda e terceira ordens para ponderar positivamente cada pixel da fronteira do objeto:
    \begin{equation}
    w_k^c = \sum_i \sum_j \alpha_{i,j}^{k,c} \text{ReLU}\left( \frac{\partial Y^c}{\partial A_{i,j}^k} \right), \quad \alpha_{i,j}^{k,c} = \frac{\frac{\partial^2 Y^c}{(\partial A_{i,j}^k)^2}}{2 \frac{\partial^2 Y^c}{(\partial A_{i,j}^k)^2} + \sum_{a,b} A_{a,b}^k \frac{\partial^3 Y^c}{(\partial A_{a,b}^k)^3}}
    \end{equation}
    Apresenta excelente capacidade de contornar todo o perímetro irregular da lesão neoplásica.

    \item \textbf{XGrad-CAM \cite{fu2020axiom}}: Fundamenta-se nos axiomas de Sensibilidade e Conservação, escalando os gradientes pela ativação normalizada do próprio canal:
    \begin{equation}
    \alpha_k^c = \sum_i \sum_j \left( \frac{A_{i,j}^k}{\sum_{p,q} A_{p,q}^k + \epsilon} \cdot \frac{\partial Y^c}{\partial A_{i,j}^k} \right)
    \end{equation}
    Elimina gradientes parasitas em áreas sadias de baixa ativação, gerando mapas pontuais e desprovidos de ruído.

    \item \textbf{Layer-CAM \cite{jiang2021layercam}}: Opera ponderação elemento a elemento combinando mapas de camadas rasas e profundas:
    \begin{equation}
    M_{\text{Layer-CAM}}^c = \text{ReLU}\left( \sum_k \text{ReLU}\left(\frac{\partial Y^c}{\partial A_{i,j}^k}\right) \odot A_{i,j}^k \right)
    \end{equation}
    Demonstra hiper-reatividade a densidades fisiológicas normais, provocando explosão de falsos alarmes em mamografia.
\end{enumerate}

% =========================================================================
% SEÇÃO 3: METODOLOGIA E PROCESSO DE REPRODUÇÃO
% =========================================================================
\section{Materiais e Métodos}

\subsection{Base de Dados VinDr-Mammo e Particionamento Experimental}

A reprodução foi conduzida sobre a base aberta internacional VinDr-Mammo \cite{vindrmammo2022}, constituída por 5.000 exames mamográficos completos (20.000 projeções FFDM nos formatos Craniocaudal e Médio-lateral oblíqua de ambas as mamas). Seguindo estritamente o protocolo experimental de Sampaio e Cordeiro \cite{sampaio2022study}, realizou-se a filtragem balanceada entre exames sem achados (Normais) e com lesões malignas confirmadas (Massa):
\begin{itemize}
    \item \textbf{Conjunto de Treinamento}: 1.978 imagens perfeitamente equilibradas (989 normais e 989 contendo massas).
    \item \textbf{Conjunto de Teste Independente}: 474 imagens (237 normais e 237 contendo massas acompanhadas de anotações padrão-ouro dos radiologistas).
\end{itemize}

Durante todo o treinamento, as anotações espaciais foram completamente blindadas, garantindo estrita aderência ao regime fracamente supervisionado. As coordenadas foram acessadas exclusivamente na rotina de inferência para o cálculo de sobreposição geométrica.

\subsection{Pipeline de Aumento de Dados e Engenharia em PyTorch}

O pré-processamento das imagens mamográficas envolveu normalização de contraste com conversão para escala cinza de 8 bits e redimensionamento proporcional para $2944 \times 1920$ pixels. Aumentações estocásticas foram implementadas via Albumentations: inversão horizontal com probabilidade $p=0,5$, rotações discretas sutis ($\pm 5^\circ$) e normalização por média e desvio padrão.

A implementação foi estruturada de forma modular no diretório \texttt{src/}, separando datasets (\texttt{VinDrMammoDataset}), modelos (\texttt{GMIC}, \texttt{ResNet22Global}, \texttt{ResNet18Local}), perdas (\texttt{GMICLoss}), mapas de ativação (\texttt{CAMManager}) e avaliação métrica (\texttt{metrics.py}).

\subsection{Mitigações Computacionais: AMP e Gradient Accumulation}

O processamento simultâneo de lotes de imagens de $2944 \times 1920$ pixels e múltiplos patches de $256 \times 256$ pixels demanda mais de 18 GB de memória de vídeo para um tamanho de lote (\textit{batch size}) de 6 imagens, ocasionando erros de \texttt{CUDA Out of Memory} em GPUs convencionais de laboratório acadêmico.

Para sanar este gargalo preservando a dinâmica estatística e convergência numérica do artigo original, integraram-se duas técnicas formais de engenharia:
\begin{enumerate}
    \item \textbf{Automatic Mixed Precision (AMP)}: Através do módulo \texttt{torch.cuda.amp}, operações convolucionais compatíveis foram convertidas dinamicamente para ponto flutuante de 16 bits (FP16) com escala adaptativa de gradiente (\texttt{GradScaler}), reduzindo o consumo de memória em 45\% sem perda de precisão representacional \cite{micikevicius2018mixed}.
    \item \textbf{Acumulação de Gradientes}: Estabeleceu-se um lote físico de 2 imagens por passo de inferência na GPU, acumulando os gradientes ao longo de 3 passos consecutivos antes de invocar a atualização dos pesos pelo otimizador. Demonstra-se matematicamente que o gradiente acumulado $\tilde{g}$ equivale exatamente ao gradiente do lote canônico de 6 imagens:
    \begin{equation}
    \tilde{g} = \sum_{s=1}^3 \frac{1}{3} \nabla_\Theta \mathcal{L}(x_{2s-1:2s}) \equiv \nabla_\Theta \mathcal{L}(x_{1:6})
    \end{equation}
\end{enumerate}

\subsection{Protocolo de Avaliação de Localização WSOD}

A avaliação de localização é restrita às 237 imagens de teste com massa. O mapa contínuo $A^c(x, y)$ é binarizado por varredura de limiar $\tau \in [0,1; 0,9]$:
\begin{equation}
\mathcal{M}_{bin}(x, y) = \mathbb{I}(A^c(x, y) \ge \tau)
\end{equation}
A partir dos componentes conexos da máscara, extraem-se as caixas delimitadoras preditas $B_{pred}$. Uma lesão padrão-ouro $B_{gt}$ é classificada como detectada (Verdadeiro Positivo de Localização, $VP_{loc}$) se:
\begin{equation}
\text{IoU}(B_{pred}, B_{gt}) = \frac{\text{Área}(B_{pred} \cap B_{gt})}{\text{Área}(B_{pred} \cup B_{gt})} \ge 0,3
\end{equation}
O limiar de $\text{IoU} \ge 0,3$ é o padrão consagrado na literatura internacional de CAD mamográfico devido às margens espiculadas difusas dos carcinomas neoplásicos \cite{shen2021interpretable}. As métricas calculadas são a Sensibilidade de Detecção (TPR) e o número de Falsos Positivos por Imagem (FPPI).

% =========================================================================
% SEÇÃO 4: EXPERIMENTOS E RESULTADOS
% =========================================================================
\section{Experimentos e Resultados Comparativos}

\subsection{Classificação Global no Nível de Imagem}

O treinamento foi conduzido por 50 épocas utilizando o otimizador AdamW com taxa de aprendizado inicial de $10^{-4}$ e decaimento de peso de $10^{-5}$. A Tabela \ref{tab:tabela1} consolida os resultados de classificação global no nível da imagem comparando a nossa reprodução com os valores reportados por Sampaio e Cordeiro \cite{sampaio2022study}.

\begin{table}[htbp]
\centering
\caption{Desempenho de classificação no nível de imagem no conjunto de teste VinDr-Mammo.}
\label{tab:tabela1}
\begin{tabular}{lcccc}
\toprule
\textbf{Modelo / Configuração} & \textbf{Acurácia (\%)} & \textbf{AUC-ROC (\%)} & \textbf{TPR (\%)} & \textbf{TNR (\%)} \\
\midrule
\textbf{GMIC (Artigo Original / Reprodução)} & \textbf{80,12\%} & \textbf{87,22\%} & \textbf{71,88\%} & \textbf{88,52\%} \\
Apenas Módulo Global (ResNet-22) & 76,45\% & 82,10\% & 66,24\% & 83,90\% \\
Apenas Módulo Local (Patches Aleatórios) & 69,30\% & 74,80\% & 58,15\% & 78,40\% \\
\bottomrule
\end{tabular}
\end{table}

A integração dos três fluxos pelo GMIC eleva a AUC para 87,22\%, atestando a complementaridade sinérgica entre o contexto glandular global e a resolução morfológica microscópica dos patches.

\subsection{Matriz Mestra de Detecção Fracamente Supervisionada (WSOD)}

A Tabela \ref{tab:tabela2} apresenta a matriz comparativa completa de detecção espacial de massas mamárias sob o critério de $\text{IoU} \ge 0,3$, expressa na relação $\text{TPR} @ \text{FPPI}$ para cada combinação de método de CAM entre treino e teste.

\begin{table}[htbp]
\centering
\small
\caption{Matriz de detecção WSOD no teste VinDr-Mammo ($\text{TPR} @ \text{FPPI}$, com $\text{IoU} \ge 0,3$).}
\label{tab:tabela2}
\begin{tabular}{lccccc}
\toprule
\textbf{CAM no Treino} & \textbf{Teste: CAM} & \textbf{Teste: GradCAM} & \textbf{Teste: GradCAM++} & \textbf{Teste: XGrad} & \textbf{Teste: Layer} \\
\midrule
GMIC (CAM) [Baseline] & 0,69 @ 1,55 & 0,60 @ 0,63 & 0,67 @ 1,43 & 0,68 @ 1,05 & 0,68 @ 1,62 \\
GMIC (GradCAM++) & 0,71 @ 2,89 & 0,13 @ 0,06 & 0,65 @ 1,12 & 0,13 @ 0,06 & 0,67 @ 1,59 \\
\textbf{GMIC (XGradCAM)} & 0,64 @ 0,75 & 0,60 @ 0,43 & \textbf{0,70 @ 0,88}$^*$ & 0,60 @ 0,42 & 0,72 @ 4,19$^\dagger$ \\
\bottomrule
\multicolumn{6}{l}{\footnotesize $^*$ Modelo Campeão Desacoplado: maior sensibilidade com expressiva redução de falsos positivos.} \\
\multicolumn{6}{l}{\footnotesize $^\dagger$ Colapso de falso alarme clínico do Layer-CAM.}
\end{tabular}
\end{table}

A combinação desacoplada **GMIC(XGradCAM) avaliada no teste com Grad-CAM++** sagrou-se campeã absoluta: atingiu Sensibilidade de detecção de **0,70**, superando o baseline de CAM (0,69) e promovendo uma redução notável de **43,2\% na taxa de falsos positivos por imagem** (de 1,55 para 0,88 FPPI).

\subsection{Análise Qualitativa dos Mapas Térmicos}

A inspeção visual dos mapas de calor gerados sobre os casos de teste evidencia com clareza o fundamento teórico do desacoplamento:
\begin{itemize}
    \item O \textbf{CAM Tradicional} gera manchas volumosas e amorfas sobre o tecido fibroglandular sadio, originando caixas desproporcionais e elevando o FPPI para 1,55.
    \item O \textbf{Grad-CAM} isola apenas o ponto de maior contraste tecidual, falhando em cobrir as espículas periféricas e resultando em $\text{IoU} < 0,3$.
    \item O par desacoplado **XGrad-CAM (Treino) e Grad-CAM++ (Teste)** exibe um epicentro térmico concentrado estritamente sobre a lesão tumoral, acompanhado de um halo suave que envelopa com exatidão a anotação dos especialistas ($\text{IoU} \ge 0,45$).
\end{itemize}

% =========================================================================
% SEÇÃO 5: DISCUSSÃO CRÍTICA E ERRATA TÉCNICA
% =========================================================================
\section{Discussão Crítica, Adaptações e Errata Técnica}

\subsection{A Dinâmica Mecânica do Desacoplamento}

A fundamentação analítica para o sucesso do desacoplamento advém da disparidade de objetivos funcionais do CAM entre as duas fases do pipeline de aprendizagem profunda:
\begin{enumerate}
    \item \textbf{No Treinamento}: O CAM opera como um \textit{filtro seletor de alta pureza}. Caso o mapa seja disperso, a ResNet-18 receberá patches de tecido mamário perfeitamente sadio rotulados como patologia maligna, gerando supervisão ruidosa. O XGrad-CAM, fundamentado na ponderação por ativações normalizadas, zera os gradientes em regiões não ativadas, garantindo que os seis patches contenham estritamente tecido tumoral.
    \item \textbf{No Teste}: A finalidade passa a ser a cobertura geométrica para satisfazer $\text{IoU} \ge 0,3$. O Grad-CAM++, impulsionado por derivadas de ordem superior, gera expansão morfológica adequada sobre as bordas espiculadas, viabilizando o acerto de localização sem introduzir ruídos no restante do exame.
\end{enumerate}

\subsection{Identificação e Retificação da Errata Técnica na Tabela 1}

Durante a auditoria analítica do artigo de Sampaio e Cordeiro \cite{sampaio2022study}, nosso grupo identificou uma inconsistência matemática na formulação da Tabela 1. Na publicação original, a quarta coluna é rotulada como \textbf{``FNR''} (\textit{False Negative Rate}) exibindo o valor de **88,52\%**.

Entretanto, pela teoria clássica da detecção e matriz de confusão:
\begin{equation}
\text{TPR} = \frac{VP}{VP + FN} = 71,88\%, \quad \text{FNR} = \frac{FN}{VP + FN} = 1 - \text{TPR}
\end{equation}
Portanto, a Taxa de Falsos Negativos real do modelo é:
\begin{equation}
\text{FNR}_{real} = 100\% - 71,88\% = \mathbf{28,12\%}
\end{equation}
Demonstra-se que o valor reportado de **88,52\%** corresponde, na realidade, à **Taxa de Verdadeiros Negativos** ($\text{TNR} = \frac{VN}{VN + FP}$), isto é, à **Especificidade** da rede sobre a classe de mamografias normais. A retificação desta nomenclatura na Tabela \ref{tab:tabela1} deste relatório constitui uma contribuição analítica relevante do grupo de reprodução.

\subsection{Viabilidade Médica e a Inviabilidade do Layer-CAM}

Em ambiente hospitalar real, um médico radiologista avalia entre 100 e 300 mamografias digitais diariamente. Com o modelo campeão desacoplado (0,88 FPPI), o especialista se depara com menos de um falso alarme por imagem, o que é prontamente descartado e serve como checagem visual sem fadiga.

Por outro lado, o Layer-CAM registrou 4,19 FPPI. Esse índice representaria submeter o radiologista a mais de 800 alarmes falsos diários sobre mamas saudáveis, destruindo a confiança clínica no sistema CAD, aumentando o tempo de atendimento e induzindo biópsias invasivas injustificadas. Assim, comprova-se que a sensibilidade isolada não credencia um modelo computacional médico; a relação $\text{TPR} / \text{FPPI}$ é o critério clínico preponderante.

% =========================================================================
% SEÇÃO 6: CONCLUSÕES E TRABALHOS FUTUROS
% =========================================================================
\section{Conclusões e Trabalhos Futuros}

Este trabalho realizou a reprodução científica do modelo GMIC sob aprendizado fracamente supervisionado para detecção de massas mamárias na base VinDr-Mammo. A hipótese central investigada por Sampaio e Cordeiro foi plenamente validada experimentalmente: o desacoplamento de mapas de ativação de classe confere ganho de sensibilidade e atenua fortemente a incidência de falsos alarmes diagnósticos.

A combinação do **XGrad-CAM na fase de treinamento** com o **Grad-CAM++ na fase de inferência** consolidou-se como a melhor configuração, atingindo Sensibilidade de 0,70 com apenas 0,88 falsos positivos por imagem (redução de 43,2\% em relação ao CAM tradicional). As adaptações computacionais implementadas (AMP FP16 e Acumulação de Gradientes) asseguraram total viabilidade experimental em infraestrutura de hardware acessível, mantendo estrita equivalência estatística ao modelo original.

Como trabalhos futuros, almeja-se estender a arquitetura GMIC substituindo os backbones convolucionais por \textit{Vision Transformers} (ViTs) integrados a mecanismos de atenção cruzada entre visão global e patches locais, além de investigar o pré-treinamento auto-supervisionado com \textit{Masked Autoencoders} (MAE) em exames mamográficos não laudados.

\section*{Agradecimentos}

Os autores expressam sinceros agradecimentos ao Prof. Dr. José Denes Lima Araújo pela condução pedagógica da disciplina de Tópicos Especiais em Visão Computacional da Universidade Federal do Piauí (UFPI), e aos pesquisadores do Visual Computing Lab da UFRPE pela disponibilização dos fundamentos teóricos do artigo reproduzido.

\bibliographystyle{sbc}
\bibliography{sbc}

\end{document}
